\documentclass[11pt,a4paper]{article}

% --- Paket LaTeX Standar ---
\usepackage[utf8]{inputenc}
\usepackage[bahasa]{babel}
\usepackage[margin=1in]{geometry}
\usepackage{amsmath,amsfonts,amssymb}
\usepackage{booktabs}
\usepackage{xcolor}
\usepackage{graphicx}
\usepackage{listings}
\usepackage{enumitem}
\usepackage{fancyhdr}
\usepackage{tcolorbox}
\usepackage{hyperref}
\usepackage{mathptmx}      % Menggunakan font Times New Roman
\usepackage{parskip}       % Spasi antar paragraf modern & tanpa indentasi paragraf awal
\renewcommand{\arraystretch}{1.3} % Membuat tinggi baris tabel lebih renggang & profesional

% --- Konfigurasi Hyperref ---
\hypersetup{
    colorlinks=true,
    linkcolor=blue,
    filecolor=magenta,      
    urlcolor=cyan,
    pdftitle={Scope of Work (SOW) - Modul Akuntansi PT Semadam},
    pdfpagemode=FullScreen,
}

% --- Konfigurasi Warna & Kotak Pesan ---
\definecolor{primarycolor}{HTML}{0f172a} % Slate 900
\definecolor{secondarycolor}{HTML}{0284c7} % Sky 600
\definecolor{alertbg}{HTML}{f8fafc}
\definecolor{alertborder}{HTML}{cbd5e1}
\definecolor{alerttext}{HTML}{334155}

\newtcolorbox{notealert}[1][Catatan Penting]{
    colback=alertbg,
    colframe=alertborder,
    coltext=alerttext,
    fonttitle=\bfseries,
    title=#1,
    arc=4px,
    boxrule=1pt,
    left=10pt, right=10pt, top=8pt, bottom=8pt
}

% --- Konfigurasi Header & Footer ---
\pagestyle{fancy}
\fancyhf{}
\fancyhead[L]{\small\color{gray}Sistem ERP PT Semadam - Scope of Work (SOW)}
\fancyhead[R]{\small\color{gray}Modul Akuntansi}
\fancyfoot[L]{\small\color{gray}\copyright~2026 PT Semadam}
\fancyfoot[R]{\small\thepage}
\renewcommand{\headrulewidth}{0.4pt}
\renewcommand{\footrulewidth}{0.4pt}

% --- INFORMASI DOKUMEN ---
\title{
    \vspace{-1.5cm}
    \color{primarycolor}\textbf{\Huge SCOPE OF WORK (SOW)}\\
    \vspace{0.3cm}
    \color{secondarycolor}\Large\textbf{Pengembangan Modul Akuntansi \& Keuangan}\\
    \vspace{0.2cm}
    \large\textbf{Sistem Informasi ERP Terintegrasi PT Semadam}
}
\author{\textbf{PT Semadam - Tim Pengembang Sistem}}
\date{29 Mei 2026}

\begin{document}

\maketitle

\thispagestyle{empty}
\begin{center}
    \vspace{1.5cm}
    \begin{tabular}{ll}
        \toprule
        \textbf{Detail Dokumen} & \textbf{Informasi} \\
        \midrule
        Nama Proyek & Sistem ERP Terintegrasi PT Semadam \\
        Modul Utama & Akuntansi (Accounting - Consolidated Ledger Hub) \\
        Versi Dokumen & 1.0.0 (Final Draft) \\
        Status Dokumen & Siap Disetujui (Ready for Signature) \\
        Bahasa & Bahasa Indonesia \\
        \bottomrule
    \end{tabular}
\end{center}

\newpage
\tableofcontents
\newpage

% ==========================================
\section{PENDAHULUAN \& LATAR BELAKANG}
% ==========================================
PT Semadam sedang melakukan modernisasi infrastruktur sistem informasinya dari sistem warisan (\textit{legacy}) berbasis Clipper/dBase (DOS) menjadi \textbf{Sistem ERP Terintegrasi modern berbasis Web}. Pengembangan ERP ini dibagi menjadi 5 modul utama yang berjalan secara paralel, di mana \textbf{Modul Akuntansi (Accounting)} bertindak sebagai \textbf{Hub Keuangan Konsolidasi (General Ledger)}.

Dokumen \textit{Scope of Work} (SOW) ini mendefinisikan batas pekerjaan, tahapan implementasi, tanggung jawab para pihak, serta kriteria keberterimaan (\textit{acceptance criteria}) khusus untuk pengembangan \textbf{Modul Akuntansi}. Dokumen ini menjadi kesepakatan formal antara Tim Pengembang dan Manajemen PT Semadam.

% ==========================================
\section{TUJUAN PROYEK}
% ==========================================
Tujuan dari pekerjaan pengembangan Modul Akuntansi ini adalah:
\begin{enumerate}
    \item \textbf{Menggantikan Sistem Clipper}: Menghilangkan ketergantungan pada sistem DOS dan database \texttt{.DBF} yang berisiko tinggi terhadap kehilangan data.
    \item \textbf{Pusat Konsolidasi (Financial Hub)}: Menyediakan \textit{endpoint} integrasi yang aman dan efisien untuk menerima jurnal otomatis/append dari 4 modul operasional (Inventory, Kas/Bank, Payroll, Asset).
    \item \textbf{Optimasi Performa Laporan}: Menerapkan caching canggih via tabel \texttt{laporan\_manajemen\_netral} (LNET) dan Server Actions untuk kalkulasi Neraca Percobaan (Trial Balance) berkecepatan tinggi.
    \item \textbf{Desain Ramah Akuntan (Accountant-Friendly UI)}: Menyediakan kalkulator audit mengambang (\textit{floating calculator}), navigasi keyboard shortcut penuh untuk input jurnal, dan tata letak laporan keuangan yang optimal untuk dicetak (\textit{print-optimized}).
\end{enumerate}

\newpage

% ==========================================
\section{RUANG LINGKUP PEKERJAAN (SCOPE OF WORK)}
% ==========================================
Pekerjaan pengembangan Modul Akuntansi ini dibagi menjadi 6 Fase Utama yang mencakup perancangan database, pengembangan logika backend, pengembangan antarmuka (frontend), integrasi modul, hingga pengujian kinerja.

\subsection{Fase 1: Perancangan \& Migrasi Database (Supabase PostgreSQL)}
\begin{itemize}
    \item \textbf{Implementasi DDL Database}: Membuat 6 tabel utama akuntansi di database Supabase sesuai dengan spesifikasi SRS:
    \begin{enumerate}
        \item \texttt{public.master\_unit} (Data Unit/Kebun)
        \item \texttt{public.master\_rekening} (Chart of Accounts - COA)
        \item \texttt{public.jurnal\_transaksi} (Ledger Transaksi)
        \item \texttt{public.saldo\_awal} (Saldo Awal per Unit/Bulan/Tahun)
        \item \texttt{public.anggaran\_rabi} (Rencana Anggaran Biaya Perkebunan)
        \item \texttt{public.laporan\_manajemen\_netral} (Cache LNET Laporan Manajemen)
    \end{enumerate}
    \item \textbf{Penerapan Kebijakan Keamanan (RLS)}: Menulis skrip Row Level Security (RLS) PostgreSQL untuk membatasi hak akses data berdasarkan autentikasi pengguna.
    \item \textbf{Optimasi Indeks Relasional}: Membuat indeks komposit pada \texttt{jurnal\_transaksi} (\texttt{idx\_jurnal\_filter} pada kolom \texttt{koke}, \texttt{kobu}, \texttt{thn}, \texttt{norek}) untuk memastikan kecepatan query pencarian data historis skala besar.
    \item \textbf{Seeding Data Awal}: Melakukan migrasi data master COA historis PT Semadam (sekitar 300+ rekening perkebunan) ke dalam tabel \texttt{master\_rekening}.
\end{itemize}

\subsection{Fase 2: Global State Context \& Infrastruktur Header (Frontend)}
\begin{itemize}
    \item \textbf{Accounting Global Context (\texttt{AccountingContext})}: Membangun React Context penyimpan \textit{state} global pembukuan aktif:
    \begin{itemize}
        \item \texttt{KOKE} (Kode Kebun Aktif)
        \item \texttt{BULAN} (Bulan Pembukuan Aktif, 01-12)
        \item \texttt{TAHUN} (Tahun Pembukuan Aktif)
    \end{itemize}
    \item \textbf{Global Selector UI}: Mengintegrasikan selektor unit (\texttt{KOKE}), bulan, dan tahun pada \texttt{SiteHeader} (Navbar utama) menggunakan komponen Radix UI / Shadcn. Setiap perubahan pada selektor ini akan secara otomatis memicu pemuatan ulang (\textit{refetching}) data pada halaman aktif tanpa reload halaman penuh.
\end{itemize}

\subsection{Fase 3: Form Input Jurnal Transaksi \& Utility Widget}
\begin{itemize}
    \item \textbf{Form Input Jurnal Dinamis}:
    \begin{itemize}
        \item Mendukung entri baris Debet dan Kredit secara dinamis (tambah/hapus baris cepat).
        \item Sistem autokomplit pencarian nomor/nama rekening dari \texttt{master\_rekening}.
        \item Validasi \textit{Real-Time Balance}: Jurnal tidak dapat disimpan jika jumlah total Debet tidak sama dengan Kredit (harus \textit{balance}).
    \end{itemize}
    \item \textbf{Navigasi Keyboard Penuh}: Mengimplementasikan \textit{keyboard listeners} (Hotkeys) untuk akuntan:
    \begin{itemize}
        \item \texttt{Ctrl + S} untuk menyimpan jurnal secara instan.
        \item \texttt{Tab} / \texttt{Shift + Tab} untuk navigasi antar field input.
        \item \texttt{Esc} untuk membatalkan entri.
    \end{itemize}
    \item \textbf{Floating Audit Calculator Widget}: Membuat widget kalkulator mengambang yang dapat dipanggil di pojok kanan bawah layar untuk membantu akuntan melakukan perhitungan cepat saldo audit tanpa membuka aplikasi eksternal.
\end{itemize}

\subsection{Fase 4: Backend Processing \& Calculation Engine (Server Actions)}
\begin{itemize}
    \item \textbf{Server Action \texttt{calculateTrialBalance}}:
    \begin{itemize}
        \item Mengimplementasikan logika \textit{Server-Side Memory Joining} antara saldo awal pembukuan (\texttt{saldo\_awal}) dengan mutasi debet/kredit yang tercatat pada \texttt{jurnal\_transaksi} untuk periode bulan dan tahun terpilih.
        \item Mengembalikan data baris Neraca Percobaan (Trial Balance) yang berisi Saldo Awal, Mutasi Debet, Mutasi Kredit, dan Saldo Akhir.
    \end{itemize}
    \item \textbf{Server Action \texttt{prosesLNET}}:
    \begin{itemize}
        \item Menghitung agregasi biaya aktual per divisi/afdeling dari tabel \texttt{jurnal\_transaksi}.
        \item Membandingkan biaya aktual dengan anggaran dari tabel \texttt{anggaran\_rabi}.
        \item Melakukan kalkulasi deviasi biaya dan menyimpan hasilnya (biaya aktual s/d bulan ini, anggaran s/d bulan ini, persentase efisiensi) ke dalam tabel cache \texttt{laporan\_manajemen\_netral}.
    \end{itemize}
\end{itemize}

\subsection{Fase 5: Modul Integrasi \& Import Utility (Append Engine)}
\begin{itemize}
    \item \textbf{Append Engine Interface}: Membuat UI khusus untuk melakukan proses \textit{append} (penggabungan data) transaksi bulanan dari file operasional eksternal (csv/xlsx) dari modul Gudang (Inventory) dan Kas/Bank.
    \item \textbf{Validasi Integritas File Append}:
    \begin{itemize}
        \item Verifikasi kolom wajib (Tanggal, No. Jurnal, No. Rekening, Deskripsi, Nominal Debet/Kredit).
        \item Verifikasi keberadaan Kode Kebun (\texttt{KOKE}) dan Nomor Perkiraan (\texttt{norek}) terhadap tabel master.
    \end{itemize}
    \item \textbf{Transaction Atomicity}: Memastikan proses penulisan ribuan baris jurnal ke tabel \texttt{jurnal\_transaksi} berjalan dalam satu transaksi database PostgreSQL (\textit{All-or-Nothing transaction block}). Jika terjadi satu baris gagal validasi, seluruh proses \textit{append} dibatalkan demi keamanan data.
\end{itemize}

\subsection{Fase 6: Pelaporan Keuangan \& Print Optimization}
\begin{itemize}
    \item \textbf{Laporan Buku Besar (Running Balance Ledger)}: Tampilan buku besar per rekening dengan kolom saldo berjalan (\textit{running balance}) yang dihitung dinamis.
    \item \textbf{Laporan Manajemen Perkebunan (LM-13 \& LNET)}: Membuat dasbor visualisasi perbandingan Realisasi vs Anggaran Biaya per Afdeling untuk melacak efisiensi biaya perawatan tanaman kelapa sawit (TBM/TM) dan biaya pengolahan pabrik.
    \item \textbf{Print-Optimized Stylesheets}: Menulis CSS \texttt{@media print} khusus untuk memastikan semua laporan keuangan bersih dari elemen navigasi web (navbar, sidebar, tombol) dan terformat rapi pada media kertas ukuran \textbf{A4 atau Folio/F4} secara vertikal (\textit{portrait}) maupun horizontal (\textit{landscape}).
\end{itemize}

\newpage

% ==========================================
\section{DAFTAR DELIVERABLES \& MILESTONES}
% ==========================================
Tabel berikut merangkum target deliverables (hasil kerja fisik), kriteria keberterimaan khusus, dan estimasi waktu penyelesaian untuk masing-masing fase pekerjaan.

\begin{table}[h!]
\centering
\small
\begin{tabular}{p{0.5cm} p{3.2cm} p{4.2cm} p{4.5cm} l}
    \toprule
    \textbf{No.} & \textbf{Fase Pekerjaan} & \textbf{Target Deliverables} & \textbf{Kriteria Keberterimaan} & \textbf{Waktu} \\
    \midrule
    \textbf{1} & Database \& Security & Skema tabel PostgreSQL aktif di Supabase, RLS teruji, Indeks aktif. & Tabel database sukses diisi data inisiasi dan memblokir akses tidak sah. & Minggu 1 \\
    \textbf{2} & Global Context & Komponen \texttt{AccountingContext} terintegrasi di Navbar. & Perubahan selektor KOKE langsung meng-update data halaman tanpa reload. & Minggu 2 \\
    \textbf{3} & Form Jurnal \& Widget & Form Input Jurnal dengan keyboard hotkeys \& widget kalkulator. & Jurnal tidak seimbang diblokir, navigasi Tab dan Ctrl+S berfungsi 100\%. & Minggu 3 \\
    \textbf{4} & Calculation Engine & Server Action \texttt{calculateTrialBalance} \& \texttt{prosesLNET}. & Kalkulasi Neraca Percobaan selesai dalam < 1.5 detik untuk data 10.000+ baris. & Minggu 4 \\
    \textbf{5} & Append Engine & UI Import CSV/XLSX \& modul logging kesalahan. & Transaksi ter-append secara atomik, file rusak ditolak tanpa merusak database. & Minggu 5 \\
    \textbf{6} & Financial Reports & Halaman Buku Besar, Neraca, LM-13 & Laporan rapi saat dicetak ke PDF/printer kertas (layout tidak terpotong). & Minggu 6 \\
    \bottomrule
\end{tabular}
\caption{Tabel Jadwal Milestones dan Kriteria Serah Terima Modul Akuntansi}
\end{table}

% ==========================================
\section{BATASAN PEKERJAAN (OUT OF SCOPE)}
% ==========================================
Pekerjaan yang didefinisikan di luar lingkup (\textit{Out of Scope}) pengerjaan modul ini meliputi:
\begin{enumerate}
    \item \textbf{Pengembangan 4 Modul Operasional}: Pembuatan internal fungsionalitas modul Inventory, Kas/Bank, Payroll, dan Asset itu sendiri (karena dikerjakan oleh tim terpisah secara paralel). Akuntansi hanya bertanggung jawab atas \textbf{penerimaan data jurnal/append} dari modul tersebut.
    \item \textbf{Integrasi Direct Bank API}: Pengambilan data mutasi rekening bank secara otomatis (\textit{Open Banking API}) tidak tercakup. Penginputan data bank menggunakan import manual file Excel/CSV pada modul Kas/Bank.
    \item \textbf{Entri Data Historis Clipper}: Pengisian data pembukuan historis PT Semadam tahun-tahun sebelumnya secara manual. Tim pengembang hanya menyediakan skrip migrasi/seeding awal, sedangkan entri operasional dilakukan oleh Tim Akuntan PT Semadam.
    \item \textbf{Pengadaan Perangkat Keras}: Pembelian komputer server, printer, atau penyediaan jaringan internet fisik di lokasi kebun PT Semadam.
\end{enumerate}

\newpage

% ==========================================
\section{KRITERIA KEBERTERIMAAN UTAMA}
% ==========================================
Modul Akuntansi ini dianggap selesai dan dapat diserahterimakan jika memenuhi kriteria berikut:
\begin{itemize}
    \item \textbf{Fungsional}: Semua fitur (Selector KOKE, Input Jurnal, Trial Balance, Append, LNET, dan LM-13) berfungsi sesuai skenario SRS.
    \item \textbf{Kecepatan Laporan}: Neraca Percobaan dan Laporan Manajemen LNET harus ter-load di bawah \textbf{1.5 detik} menggunakan mekanisme cache dan Server Actions.
    \item \textbf{Keamanan RLS}: Uji coba penetrasi membuktikan bahwa pengguna dari Kebun A (\texttt{KOKE: 01}) tidak dapat melihat atau memodifikasi data milik Kebun B (\texttt{KOKE: 02}).
    \item \textbf{Kemudahan Cetak}: Pengujian cetak fisik atau cetak ke berkas PDF pada peramban menunjukkan dokumen terformat rapi pada ukuran A4/F4 tanpa terpotong (\textit{overflow}).
    \item \textbf{Akurasi Perhitungan}: Hasil kalkulasi Debet/Kredit pada Neraca Percobaan 100\% konsisten dengan akumulasi saldo transaksi.
\end{itemize}

% ==========================================
\section{STRUKTUR TIM \& KOORDINASI}
% ==========================================
\begin{itemize}
    \item \textbf{Project Sponsor}: Direktur Keuangan PT Semadam.
    \item \textbf{Product Owner}: Kepala Akuntan PT Semadam (menyediakan spesifikasi akun \& alur audit).
    \item \textbf{Lead Developer}: Arsitek Next.js \& Database Supabase.
    \item \textbf{System Integrator}: Penanggung jawab integrasi modul Kas/Bank \& Inventory.
\end{itemize}

\begin{notealert}[Sistem Koordinasi]
Koordinasi pengerjaan dilakukan melalui \textit{Daily Standup} mingguan secara daring dan peninjauan berkala hasil build pada server staging setiap akhir fase.
\end{notealert}

\vspace{2cm}

\begin{center}
\textit{Disetujui oleh,}
\end{center}

\vspace{1cm}

\begin{table}[h!]
\centering
\begin{tabular}{cc}
    \textbf{Untuk PT Semadam} & \textbf{Untuk Tim Pengembang} \\
    \vspace{2cm} & \vspace{2cm} \\
    \rule{5cm}{0.5pt} & \rule{5cm}{0.5pt} \\
    Manajemen PT Semadam & Pimpinan Tim Pengembang \\
\end{tabular}
\end{table}

\end{document}
